% =========================================================
%  config/packages.tex
%  Seluruh paket LaTeX untuk Buku Ajar Pengantar Pemrograman Game
%  Aturan: paket global hanya didefinisikan di berkas ini,
%  tidak di chapter.tex atau section-*.tex (Arsitektur §36, Rule 3).
% =========================================================

% --- Pengkodean & bahasa --------------------------------
\usepackage[utf8]{inputenc}
\usepackage[T1]{fontenc}
\usepackage{lmodern}
\usepackage[english,indonesian]{babel}

% --- Tata letak halaman ---------------------------------
\usepackage[
  a4paper,
  left=2.9cm, right=2.4cm, top=2.8cm, bottom=2.8cm,
  headsep=0.5cm, footskip=1.0cm
]{geometry}
\usepackage{pdflscape}
\usepackage{fancyhdr}
\usepackage{ragged2e}
\usepackage{microtype}

% --- Tabel ----------------------------------------------
% Catatan: opsi [table] pada xcolor (colortbl) tidak dipakai karena
% colortbl pada instalasi ini menuntut array.sty yang lebih baru
% (\insert@pcolumn tidak terdefinisi). Pewarnaan baris tabel digantikan
% oleh penebalan garis header pada config/styles.tex.
\usepackage{booktabs}
\usepackage{array}
\usepackage{tabularx}
\usepackage{longtable}
\usepackage{multirow}
\usepackage{makecell}
\usepackage{xcolor}

% --- Gambar & posisi mengambang -------------------------
\usepackage{graphicx}
\usepackage{float}

% --- Tipografi & daftar ---------------------------------
\usepackage{enumitem}
\usepackage{titlesec}
\usepackage{tocloft}
\usepackage{caption}
\usepackage{needspace}
\usepackage{etoolbox}

% --- Matematika -----------------------------------------
\usepackage{amsmath}
\usepackage{amssymb}
\usepackage{amsthm}
\usepackage{mathtools}

% --- Kode program ---------------------------------------
\usepackage{listings}
\usepackage{upquote}

% --- Diagram --------------------------------------------
\usepackage{tikz}
\usetikzlibrary{
  arrows.meta,
  positioning,
  shapes.geometric,
  shapes.misc,
  calc,
  fit,
  backgrounds,
  matrix,
  chains,
  decorations.pathreplacing,
  patterns
}
\usepackage{pgfplots}
\pgfplotsset{compat=1.18}

% --- Kotak informasi ------------------------------------
\usepackage{tcolorbox}

% --- Kompilasi modular ----------------------------------
% subfiles memuat paket import: ketika \subfile{chapters/XX/chapter}
% dijalankan, \input@path dan \Ginput@path disetel ke direktori bab,
% sehingga \includegraphics{figures/...}, \input{contoh}, dan
% \subfile{section-01} tetap benar pada kompilasi bab mandiri maupun
% pada kompilasi seluruh buku.
\usepackage{subfiles}

% --- Rujukan silang antar-bab untuk kompilasi bab mandiri
% (Arsitektur §40). Label dari main.aux diimpor bila berkasnya ada.
% Dipakai xr-hyper (bukan xr) karena hyperref sudah dimuat; xr biasa
% menimbulkan galat "\XR@prefix undefined" pada \cite.
\usepackage{xr-hyper}

% --- Tautan (dimuat paling akhir) -----------------------
\usepackage{url}
\usepackage{xurl}
\usepackage[
  bookmarks=true,
  bookmarksnumbered=true,
  bookmarksopen=true,
  breaklinks=true
]{hyperref}
